% -----------------------------------------------------------
\section{Remission, Remit}
% -----------------------------------------------------------
	\label{REMISSION}
	
% % ----------------------
% \subsection{See also}
% % ----------------------

% ----------------------
\subsection{1828 American Dictionary of the English Language} % *1828
% ----------------------
	\label{REMISSION-1828}
	
	% -----
	\subsubsection{Remission}
	% -----

		\begin{compactitem}
			\item[1] Abatement; relaxation; moderation; as the \textit{remission} of extreme rigor.
			\item[2] Abatement; diminution of intensity; as the \textit{remission} of the sun's heat; the \textit{remission} of cold; the \textit{remission} of close study or of labor.
			\item[3] Release; discharge or relinquishment of a claim or right; as the \textit{remission} of a tax or duty.
			\item[4] In medicine, abatement; a temporary subsidence of the force or violence of a disease or of pain, as distinguished from intermission, in which the disease leaves the patient entirely for a time.
			\item[5] Forgiveness; pardon; that is, the giving up of the punishment due to a crime; as the \textit{remission} of sins.
			\item[6] The act of sending back. [Not in use.]
		\end{compactitem}
		
	% -----
	\subsubsection{Remit (transitive)}
	% -----
	
		\begin{compactitem}
			\item[1] To relax, as intensity; to make less tense or violent.
			\item[2] To forgive; to surrender the right of punishing a crime; as, to remit punishment.
			\item[3] To pardon, as a fault or crime.
			\item[4] To give up; to resign.
			\item[5] To refer; as a clause that remitted all to the bishop's discretion.
			\item[6] To send back.
			\item[7] To transmit money, bills or other thing in payment for goods received. American merchants remit money, bills of exchange or some species of stock, in payment for British goods.
			\item[8] To restore.
		\end{compactitem}

% % ----------------------
% \subsection{Oxford English Dictionary} % *OED
% % ----------------------
	% \label{REMISSION-OED}

	% \begin{compactitem}
		% \item[]
	% \end{compactitem}

% ----------------------
\subsection{Scriptural Usage} % *SCRIPTURES
% ----------------------
	\label{REMISSION-SCRIPTURES}

	% % -----
	% \subsubsection{Old Testament} % *OT
	% % -----
		% \label{REMISSION-OT}
	
		% % --
		% \paragraph{KJV}
		% % --
			% \label{REMISSION-OT-KJV}
	
			% \begin{tabular}{ll}
				% Total occurrences in the OT & 
			% \end{tabular}
			
			% \begin{compactdesc}
				% \item[]
			% \end{compactdesc}
			
		% % --
		% \paragraph{Hebrew Bible}
		% % --
			% \label{REMISSION-HEBREW}
		
			% %
			% \subparagraph{\texorpdfstring{\he{sample word}}{}}
			% %
				% \label{PAGA} % whatever is in step bible without periods
			
				% \begin{compactdesc}
					% \item[HALOT , PDF ] \textbf{}
						% \begin{compactitem}
							% \item[1]
						% \end{compactitem}
					% \item[BDB , PDF ] \textbf{}
						% \begin{compactitem}
							% \item[1]
						% \end{compactitem}
					% \item[DCH PDF ] \textbf{}
						% \begin{compactitem}
							% \item[1]
						% \end{compactitem}
				% \end{compactdesc}
				
				% \begin{compactdesc}
					% \item[Some OT uses] \textbf{}
						% \begin{compactdesc}
							% \item[]
						% \end{compactdesc}
				% \end{compactdesc}
			
		% % --
		% \paragraph{Septuagint}
		% % --
			% \label{REMISSION-LXX}
		
			% %
			% \subparagraph{\texorpdfstring{\gr{sample word}}{}}
			% %
		
	% -----
	\subsubsection{New Testament} % *NT
	% -----
		\label{REMISSION-NT}
	
		% --
		\paragraph{KJV}
		% --
			\label{REMISSION-NT-KJV}
			
	
			\begin{tabular}{ll}
				Total occurrences in the NT & 12
			\end{tabular}
			
			\begin{compactdesc}
				\item[Matthew 26:28] For this is my blood of the new testament, which is shed for many for the remission of sins.
				\item[Mark 1:4] John did baptize in the wilderness, and preach the baptism of repentance for the remission of sins.
				\item[Luke 1:77] To give knowledge of salvation unto his people by the remission of their sins,
				\item[Luke 3:3] And he came into all the country about Jordan, preaching the baptism of repentance for the remission of sins;
				\item[Luke 24:47] And that repentance and remission of sins should be preached in his name among all nations, beginning at Jerusalem.
				\item[John 20:23] Whose soever sins ye remit, they are remitted unto them; and whose soever sins ye retain, they are retained.
				\item[Acts 2:38] Then Peter said unto them, Repent, and be baptized every one of you in the name of Jesus Christ for the remission of sins, and ye shall receive the gift of the Holy Ghost.
				\item[Acts 10:43] To him give all the prophets witness, that through his name whosoever believeth in him shall receive remission of sins.
				\item[Romans 3:25] Whom God hath set forth to be a propitiation through faith in his blood, to declare his righteousness for the remission of sins that are past, through the forbearance of God;
				\item[Hebrews 9:22] And almost all things are by the law purged with blood; and without shedding of blood is no remission.
				\item[Hebrews 10:18] Now where remission of these is, there is no more offering for sin.
			\end{compactdesc}
			
		% --
		\paragraph{Greek New Testament} % *GREEK
		% --
			\label{REMISSION-GREEK}
		
			%
			\subparagraph{\texorpdfstring{\gr{ἄφεσις}}{}}
			%
				\label{APHESIS}
			
				\begin{compactdesc}
					\item[BDAG 136a] \textbf{}
						\begin{compactitem}
							\item[1] the act of freeing and liberating from something that confines, \textit{release}
							\item[2] the act of freeing from an obligation, guilt, or punishment, \textit{pardon, cancellation}
						\end{compactitem}
					\item[LSJ 288a, PDF 361] \textbf{}
						\begin{compactitem}
							\item[1] \textit{letting go, release}
							\item[2] \textit{quittance from}
							\item[2b] \textit{forgiveness} (cites the gospels of Mark and Matthew)
							\item[3] \textit{relaxation, exhaustion}
							\item[4] \textit{divorce}
							\item[5] \textit{starting} of horses in a race
							\item[6] \textit{discharge, emission}
							\item[7] (alternate form in Aristotle)
							\item[8] \textit{release} (of water)
							\item[9] \textit{reckoning of the vital quadrant} (astrological meaning)
						\end{compactitem}
				\end{compactdesc}
				
				\begin{compactdesc}
					\item[Some NT uses] \textbf{}
						\begin{compactdesc}
							\item[Matthew 26:28] \gr{τοῦτο γάρ ἐστιν τὸ αἷμά μου τῆς διαθήκης τὸ περὶ πολλῶν ἐκχυννόμενον εἰς ἄφεσιν ἁμαρτιῶν·}
							\item[Mark 1:4] \gr{ἐγένετο Ἰωάννης ὁ βαπτίζων ἐν τῇ ἐρήμῳ κηρύσσων βάπτισμα μετανοίας εἰς ἄφεσιν ἁμαρτιῶν.}
							\item[Luke 1:77] \gr{τοῦ δοῦναι γνῶσιν σωτηρίας τῷ λαῷ αὐτοῦ ἐν ἀφέσει ἁμαρτιῶν αὐτῶν,}
							\item[Luke 3:3] \gr{καὶ ἦλθεν εἰς πᾶσαν περίχωρον τοῦ Ἰορδάνου κηρύσσων βάπτισμα μετανοίας εἰς ἄφεσιν ἁμαρτιῶν,}
							\item[Luke 24:47] \gr{καὶ κηρυχθῆναι ἐπὶ τῷ ὀνόματι αὐτοῦ μετάνοιαν καὶ ἄφεσιν ἁμαρτιῶν εἰς πάντα τὰ ἔθνη--- ἀρξάμενοι ἀπὸ Ἰερουσαλήμ·}
							\item[Acts 2:38] \gr{Πέτρος δὲ πρὸς αὐτούς· Μετανοήσατε, καὶ βαπτισθήτω ἕκαστος ὑμῶν ἐπὶ τῷ ὀνόματι Ἰησοῦ Χριστοῦ εἰς ἄφεσιν τῶν ἁμαρτιῶν ὑμῶν, καὶ λήμψεσθε τὴν δωρεὰν τοῦ ἁγίου πνεύματος·}
							\item[Acts 10:43] \gr{τούτῳ πάντες οἱ προφῆται μαρτυροῦσιν, ἄφεσιν ἁμαρτιῶν λαβεῖν διὰ τοῦ ὀνόματος αὐτοῦ πάντα τὸν πιστεύοντα εἰς αὐτόν.}
							\item[Hebrews 9:22] \gr{καὶ σχεδὸν ἐν αἵματι πάντα καθαρίζεται κατὰ τὸν νόμον, καὶ χωρὶς αἱματεκχυσίας οὐ γίνεται ἄφεσις.}
							\item[Hebrews 10:18] \gr{ὅπου δὲ ἄφεσις τούτων, οὐκέτι προσφορὰ περὶ ἁμαρτίας.}
						\end{compactdesc}
				\end{compactdesc}
				
			%
			\subparagraph{\texorpdfstring{\gr{ἀφίημι}}{}}
			%
				\label{APHIEMI}
			
				\begin{compactdesc}
					\item[BDAG 137a] \textbf{}
						\begin{compactitem}
							\item[1] to dismiss or release someone or something from a place or one's presence
							\item[2] to release from legal or moral obligation or consequence, \textit{cancel, remit, pardon}
							\item[3] to move away, with implication of causing a separation, \textit{leave, depart from}
							\item[4] to have something continue or remain in a place, \textit{leave standing/lying}
							\item[5] leave it to someone to do something, \textit{let, let go, allow, tolerate}
						\end{compactitem}
					\item[LSJ 289b, PDF 362] \textbf{}
						\begin{compactitem}
							\item[I.1] \textit{send forth, discharge}
							\item[I.2] \textit{let fall}
							\item[I.3] \textit{give up} or \textit{hand over}
							\item[II.1] \textit{send away} (of persons)
							\item[II.1b] \textit{let go, loose, set free}...in a legal sense, \textit{acquit} of a charge or engagement
							\item[II.2] \textit{send away} (of things)
							\item[III.1] \textit{leave alone}
							\item[III.2] \textit{let} (something happen)
							\item[IV.1] \textit{suffer, permit} someone to do something
							\item[V.1] \textit{break up, march, sail}
							\item[V.2] \textit{give up doing}
						\end{compactitem}
				\end{compactdesc}
				
				\begin{compactdesc}
					\item[Some NT uses] \textbf{}
						\begin{compactdesc}
							\item[John 20:23] \gr{ἄν τινων ἀφῆτε τὰς ἁμαρτίας ἀφέωνται αὐτοῖς· ἄν τινων κρατῆτε κεκράτηνται.}
						\end{compactdesc}
				\end{compactdesc}
				
			%
			\subparagraph{\texorpdfstring{\gr{πάρεσις}}{}}
			%
				\label{PARESIS}
			
				\begin{compactdesc}
					\item[BDAG 689b] \textbf{}
						\begin{compactitem}
							\item deliberate disregard, \textit{passing over, letting go unpunished} (BDAG at least lists the use in Romans 3:25 as the only occurrence of this word in the NT, and adds the following: ``The verb [\gr{παρίημι} I think] is also used of `remitting' debts and other obligations...but in Romans 3:25 Paul appears to argue that God's apparent winking at sin puts God's own uprightness at risk.'')
						\end{compactitem}
					\item[LSJ 1337b, PDF 1408] \textbf{}
						\begin{compactitem}
							\item[I] \textit{letting go, dismissal}
							\item[II.1] \textit{slackening} of strength, \textit{paralysis}
							\item[II.2] a metaphorical sense of II.1
							\item[III] \textit{remission} of debts (or sins; cites Romans 3:25)
							\item[IV] \textit{neglect}
						\end{compactitem}
				\end{compactdesc}
				
				\begin{compactdesc}
					\item[Some NT uses] \textbf{}
						\begin{compactdesc}
							\item[Romans 3:25] \gr{ὃν προέθετο ὁ θεὸς ἱλαστήριον διὰ πίστεως ἐν τῷ αὐτοῦ αἵματι εἰς ἔνδειξιν τῆς δικαιοσύνης αὐτοῦ διὰ τὴν πάρεσιν τῶν προγεγονότων ἁμαρτημάτων}
						\end{compactdesc}
				\end{compactdesc}
		
	% -----
	\subsubsection{Book of Mormon} % *BOM
	% -----
		\label{REMISSION-BOM}
	
		\begin{tabular}{ll}
			Total occurrences in the BOM & 28
		\end{tabular}
		
		\begin{compactdesc}
			\item[2 Nephi 25:26] And we talk of Christ, we rejoice in Christ, we preach of Christ, we prophesy of Christ; and we write according to our prophecies that our children may know to what source they may look for a remission of their sins. 
			\item[2 Nephi 31:17] Wherefore do the things which I have told you that I have seen, that your Lord and your Redeemer should do. For for this cause have they been shewn unto me, that ye might know the gate by which ye should enter. For the gate by which ye should enter is repentance and baptism by water, and then cometh a remission of your sins by fire and by the Holy Ghost.
			\item[{\hyperref[MOSIAH-4-2]{Mosiah 4:2--3}}]
			\item[{\hyperref[MOSIAH-4-11]{Mosiah 4:11--12}}]
			\item[{\hyperref[MOSIAH-4-20]{Mosiah 4:20}}]
			\item[{\hyperref[MOSIAH-4-26]{Mosiah 4:26}}]
			\item[{\hyperref[ALMA-4-11]{Alma 4:11--14}}]
			\item[{\hyperref[ALMA-12-34]{Alma 12:34}}]
			\item[{\hyperref[ALMA-13-16]{Alma 13:16}}]
			\item[{\hyperref[ALMA-38-8]{Alma 38:8}}]
			\item[3 Nephi 12:2]  And again, more blessed are they which shall believe in your words because that ye shall testify that ye have seen me and that ye know that I am. Yea, blessed are they which shall believe in your words and come down into the depths of humility and be baptized, for they shall be visited with fire and with the Holy Ghost and shall receive a remission of their sins. 
			\item[Moroni 8:25--26] And the firstfruits of repentance is baptism. And baptism cometh by faith unto the fulfilling the commandments; and the fulfilling the commandments bringeth remission of sins; 26 and the remission of sins bringeth meekness and lowliness of heart. And because of meekness and lowliness of heart cometh the visitation of the Holy Ghost, which Comforter filleth with hope and perfect love, which love endureth by diligence unto prayer until the end shall come, when all the saints shall dwell with God. 
		\end{compactdesc}
		
	% -----
	\subsubsection{Doctrine and Covenants} % *DC
	% -----
		\label{REMISSION-DC}
	
		\begin{tabular}{ll}
			Total occurrences in the DC & 21
		\end{tabular}
		
		\begin{compactdesc}
			\item[DC 20:37] \textit{And again, by way of commandment to the church concerning the manner of baptism}---All those who humble themselves before God, and desire to be baptized, and come forth with broken hearts and contrite spirits, and witness before the church that they have truly repented of all their sins, and are willing to take upon them the name of Jesus Christ, having a determination to serve him to the end, and truly manifest by their works that they have received of the Spirit of Christ unto the remission of their sins, shall be received by baptism into his church.
			\item[DC 21:9] For, behold, I will bless all those who labor in my vineyard with a mighty blessing, and they shall believe on his words, which are given him through me by the Comforter, which manifesteth that Jesus was crucified by sinful men for the sins of the world, yea, for the remission of sins unto the contrite heart.
			\item[DC 27:2] For, behold, I say unto you, that it mattereth not what ye shall eat or what ye shall drink when ye partake of the sacrament, if it so be that ye do it with an eye single to my glory---remembering unto the Father my body which was laid down for you, and my blood which was shed for the remission of your sins.
			\item[{\hyperref[DC55-1]{DC 55:1}}] 
			\item[{\hyperref[DC56-14]{DC 56:14--15}}] 
				\begin{compactitem}
					\item Notice especially the reason why they are not ``pardoned'' here---they ``seek to counsel in [their] own ways.'' They weren't willing to accept the Lord's counsel about certain matters.
				\end{compactitem}
			\item[{\hyperref[DC64-3]{DC 64:3--13}}]
				\begin{compactitem}
					\item These verses are important for the idea of a remission as well, but they're absolutely crucial for \hyperref[FORGIVENESS]{Forgiveness}.
				\end{compactitem}
			\item[DC 68:27] 
			\item[DC 84:64] 
			\item[DC 132:46] 
		\end{compactdesc}
		
	% -----
	\subsubsection{Pearl of Great Price} % *PGP
	% -----
		\label{REMISSION-PGP}
	
		\begin{tabular}{ll}
			Total occurrences in the PGP & 3
		\end{tabular}
		
		Two uses in the JS-H, and one in the AoF.
		
		% \begin{compactdesc}
			% \item[]
		% \end{compactdesc}
		
% % ----------------------
% \subsection{Key Phrases} % *KEYPHRASES *PHRASES
% % ----------------------
	% \label{REMISSION-PHRASES}

	% \begin{compactdesc}
		% \item[]
	% \end{compactdesc}
	
% % ----------------------
% \subsection{Bible Dictionaries} % *DICTIONARIES
% % ----------------------
	% \label{REMISSION-BD}

% % ----------------------
% \subsection{Restoration Usage} % *RESTORATION
% % ----------------------
	% \label{REMISSION-RESTORATION}

	% % -----
	% \subsubsection{Joseph Smith} % *JS
	% % -----
		% \label{REMISSION-JS}
	
		% \begin{compactdesc}
			% \item[DATE/SOURCE]
		% \end{compactdesc}
	
	% % -----
	% \subsubsection{Brigham Young} % *BY
	% % -----
		% \label{REMISSION-BY}
	
		% \begin{compactdesc}
			% \item[DATE/SOURCE]
		% \end{compactdesc}
	
% ----------------------
\subsection{Commentary and Studies} % *COMMENTARY *STUDIES
% ----------------------
	\label{REMISSION-COMMENTARY}

	% -----
	\subsubsection{Key Points}
	% -----
		\label{REMISSION-KEYS}
	
		\begin{compactitem}
			\item Receiving a remission of sin, or having your sins remitted, is to be free from the legal consequences of your actions; you are no longer liable for the punishment, whatever it is, attached to the law you violated.
			\item Remission of sin and forgiveness of sin are the same thing, and the different words just call attention to different aspects of the process; see \hyperref[REMISSION-forgiveness]{this study} below.
			\item The remission of sin comes from fufilling the commandments (sort of \hyperref[MATTHEW-19-17]{Matthew 19:17} [see the Greek in particular], \hyperref[LUKE-7-37]{Luke 7:37--49} [Christ says that the woman's sins are forgiven in response to her actions], sort of \hyperref[ALMA-12-34]{Alma 12:34}, \hyperref[MORONI-8-25]{Moroni 8:25}, \hyperref[DC1-32]{DC 1:32}, \hyperref[DC4-2]{DC 4:2} (being ``blameless''), \hyperref[DC20-37]{DC 20:37}, \hyperref[DC39-10]{DC 39:10} (``wash away your sins'' after baptism), \hyperref[DC61-33]{DC 61:33--34}, \hyperref[DC76-52]{DC 76:52}, the \hyperref[TEMPLE-ENDOWMENT-ENDOW-PRE-1990-INTRO]{introduction} to the endowment (text in the next bullet); see also \hyperref[REMISSION-commandments]{this study} below). This is a general principle, but it is not exceptionless.
				\begin{compactitem}
					\item Text from introduction to the endowment: ``you have been washed and pronounced clean, or that through your faithfulness you may become clean, from the blood and sins of this generation.''
				\end{compactitem}
			\item We retain a remission of sin by keeping in mind our relationship to God, and then acting in certain ways in response (see \hyperref[MOSIAH-4-11]{Mosiah 4:11--12}).
			\item There is an internal/mental component to receiving a remission---see \hyperref[DC55-1]{DC 55:1}.
		\end{compactitem}
		
	% -----
	\subsubsection{Remission of sins and forgiveness of sins}
	% -----
		\label{REMISSION-forgiveness}
		
		What is the relationship between the concept of ``remission'' and that of ``forgiveness''? I believe they are the same, but the two words serve to emphasize different aspects of the same process or experience.
		
		In essence, to receive a remission of sin and a forgiveness of sin is the same thing. It is to be free from the legal consequences of your actions; you are no longer liable for the punishment, whatever it is, attached to the law you violated.
		
		The word ``remission'' focuses on that legal sense and calls attention to the broken law, the consequences for breaking it, and the fact that those consequences have been officially dissolved or made void. The word ``remit,'' the verb form, has an even stronger legal sense I think.
		
		The word ``forgiveness'' focuses on the personal interaction between the person who has broken the law, and the lawgiver (or the one who can dissolve the consequences). There are many instances in the DC, for example, where the Lord includes in a revelation the point that someone's sins have been forgiven. This point often appears at the beginning of a revelation, almost as a greeting, in a way.
		
		The evidence for the equivalence between those two things is in the pattern of uses throughout the scriptures, but there are some specific things you can point to. We'll begin with \hyperref[ENOS-1-2]{Enos 1:2--6}:

			\begin{quote}
				and I will tell you of the wrestle which I had before God before that I received a remission of my sins. 3 Behold, I went to hunt beasts in the forest, and I remembered the words which I had often heard my father speak concerning eternal life and the joy of the saints; and the words of my father sunk deep into my heart, 4 and my soul hungered, and I kneeled down before my Maker, and I cried unto him in mighty prayer and supplication for mine own soul. And all the day long did I cry unto him; yea, and when the night came, I did still raise my voice high, that it reached the heavens. 5 And there came a voice unto me, saying: Enos, thy sins are forgiven thee, and thou shalt be blessed. 6 And I Enos knew that God could not lie; wherefore my guilt was swept away.
			\end{quote}
			
		Here Enos says that he had a wrestle before he obtained a ``remission'' of his sins. When the remission actually arrives, in verse 5, the Lord says, ``Enos, thy sins are forgiven thee.''
		
		Another good example is from \hyperref[MOSIAH-4-3]{Mosiah 4:2--3, 10--12, 20}:
		
			\begin{quote}
				2 And they had viewed themselves in their own carnal state, even less than the dust of the earth, and they all cried aloud with one voice, saying: O have mercy and apply the atoning blood of Christ that we may receive forgiveness of our sins and our hearts may be purified, for we believe in Jesus Christ the Son of God, who created heaven and earth and all things, who shall come down among the children of men. 3 And it came to pass that after they had spoken these words, the Spirit of the Lord came upon them, and they were filled with joy, having received a remission of their sins and having peace of conscience because of the exceeding faith which they had in Jesus Christ, which should come, according to the words which king Benjamin had spoken unto them.
				
				10 And again, believe that ye must repent of your sins and forsake them and humble yourselves before God and ask in sincerity of heart that he would forgive you. And now if you believe all these things, see that ye do them. 11 And again I say unto you, as I have said before, that as ye have come to the knowledge of the glory of God or if ye have known of his goodness and have tasted of his love and have received a remission of your sins, which causeth such exceeding great joy in your souls, even so I would that ye should remember and always retain in remembrance the greatness of God and your own nothingness and his goodness and long-suffering towards you unworthy creatures, and humble yourselves even in the depths of humility, calling on the name of the Lord daily and standing steadfastly in the faith of that which is to come, which was spoken by the mouth of the angel. 12 And behold, I say unto you that if ye do this, ye shall always rejoice and be filled with the love of God and always retain a remission of your sins; and ye shall grow in the knowledge of the glory of him that created you, or in the knowledge of that which is just and true. 
				
				20 And behold, even at this time ye have been calling on his name and begging for a remission of your sins. And hath he suffered that ye have begged in vain? Nay, he hath poured out his Spirit upon you and hath caused that your hearts should be filled with joy and hath caused that your mouths should be stopped, that ye could not find utterance, so exceeding great was your joy. 
			\end{quote}
			
		The people ask for forgiveness (verse 2), then are visited by the Spirit and receive a remission of sins (verse 3). Benjamin then instructs them in how to retain that remission (verses 10--11), and later notes that when the people were asking for forgiveness, they were ``begging for a remission'' of sins (verse 20).
		
	% -----
	\subsubsection{Remission of sins and fulfilling the commandments}
	% -----
		\label{REMISSION-commandments}
		
		There are a number of verses which explicitly state that remission of sin (or forgiveness) comes from obeying the commandments. Here are some, beginning with Moroni 8:25:
		
			\begin{quote}
				And the firstfruits of repentance is baptism. And baptism cometh by faith unto the fulfilling the commandments; and the fulfilling the commandments bringeth remission of sins;
			\end{quote}
			
		DC 1:32:
		
			\begin{quote}
				Nevertheless, he that repents and does the commandments of the Lord shall be forgiven;
			\end{quote}
			
		Many verses talk about how baptism is ``for'' or ``unto'' remission of sin; one of the most explicit, which also mentions fulfilling the commandments, is Moroni 8:11: 
		
			\begin{quote}
				And their little children need no repentance, neither baptism. Behold, baptism is unto repentance, to the fulfilling the commandments unto the remission of sins.
			\end{quote}
			
		Other verses that are useful here include \hyperref[DC55-1]{DC 55:1} and \hyperref[DC108-1]{DC 108:1}.
		
		For me these verses raise a number of interesting questions. Two I want to discuss here are:
		
			\begin{compactenum}
				\item What is it about fulfilling the commandments which brings a remission? 
				\item Which commandments are referred to in this principle?
			\end{compactenum}
			
			% --
			\paragraph{1. Why fulfilling the commandments?}
			% --
				\label{REMISSION-commandments-why}
			
				What is it about fulfilling the commandments which brings a remission?
				
				One way to think about this question is to consider what \textit{else} it could be that could bring a remission. I can see at least a couple options. One is that a remission could come merely from asking for it. The other is that a remission could come whether we want it or not---it just comes to everybody.
				
				The first of these options actually isn't a bad one and I think there are instances in the scriptures where it does actually occur. The most prominent, at least in what I've been reading lately, is in \hyperref[MOSIAH-4]{Mosiah 4}, where the people are listening to Benjamin and then ask the Lord for forgiveness. They receive a remission even though it isn't clear exactly what ``commandment'' they have obeyed in order to receive it; it seems all they've done is just ask.
				
				There's no reason to think that a remission cannot come just by asking in some circumstances. The obedience-remission principle doesn't have to be exceptionless; it could just be a general truth which holds for most cases and most times, and it would still be important to know, though there would then be other ways of receiving a remission.
				
				The second option, that a remission comes to everyone whether they want it or not, isn't as appealing, at least to me. Some people won't want it, and the scriptures make clear, especially in \hyperref[DC88]{DC 88}, that people will be given what they want, and won't be forced to accept any level of salvation that they aren't interested in.
				
				With those options out of the way, what then is it about fulfilling the commandments which brings a remission? I have been able to come up with two possible explanations, which may both work together (and maybe with others I haven't thought of) to answer this question.
				
				The first explanation is that this is just the the spiritual laws are set up. We know there are laws governing the relationship between human behavior and spiritual consequences. Some of these laws, as we see in the general obedience-remission principle, govern the reception of a remission of sin. It may be that the laws simply specify that \textit{actions} are required on our part, rather than something else, in order to activate the (good) consequences of the law.
				
				This is a good explanation but perhaps not as deep as we want it to go. In particular it raises the question of \textit{why} the laws would be set up this way, and not in some other way.
				
				The second possible explanation may help here. The gospel is about \textit{becoming} something---it is about \textit{transforming} ourselves into something new, a new creature, which is eligible for experiences that the old one could not have had. For human beings, with the way our physiology works, becoming something usually requires action over a time period, medium or long. We become something by \textit{doing} something---we become generous by doing generous actions, we come honest by doing honest actions, and so on. We become spiritual by doing spiritual actions, we become sanctified by doing sanctified actions. Aristotle taught this very clearly and he was right about it.
				
				According to this explanation, the reason why obedience is required for a remission is that obedience is action, and receiving a remission requires actions, not just holding particular beliefs. We must \textit{act} on those beliefs, not just hold them or mentally affirm them, and obedience is action.
				
				This second explanation may give us a way to understand why the laws governing the reception of a remission involve action. The laws are built for humans, to be useful for us and to give us the best shot possible; it isn't the reverse, that we must conform ourselves to laws that weren't necessarily meant for humans. And so the laws involve action, because change and transformation as humans requires action.
				
				None of this is to deny the role of God's power in our transformation. We are capable of changing our behavior, but being sanctified and becoming holy requires God's power and God's actions, not only ours.
				
			% --
			\paragraph{2. Which commandments?}
			% --
				\label{REMISSION-commandments-which}
			
				In a number of places, including multiple verses in Moroni 8, the principle of remission and commandments occurs alongside a discussion of baptism. There are also many places in the scriptures, including the NT, BOM, and DC, which mention remission in conjunction with baptism. Here are some of those verses. The first is Luke 3:3, talking about Christ:
				
					\begin{quote}
						And he came into all the country about Jordan, preaching the baptism of repentance for the remission of sins;
					\end{quote}
					
				2 Nephi 31:17:

					\begin{quote}
						Wherefore do the things which I have told you that I have seen, that your Lord and your Redeemer should do. For for this cause have they been shewn unto me, that ye might know the gate by which ye should enter. For the gate by which ye should enter is repentance and baptism by water, and then cometh a remission of your sins by fire and by the Holy Ghost. 
					\end{quote}
					
				DC 33:11:
				
					\begin{quote}
						Yea, repent and be baptized, every one of you, for a remission of your sins; yea, be baptized even by water, and then cometh the baptism of fire and of the Holy Ghost.
					\end{quote}
					
				So, while the principle as stated in DC 1:32 is given generally, elsewhere it appears together with baptism, the context making it fairly clear that the commandment referred to is baptism in particular.
				
				It seems that baptism happens to be a very special commandment. It is the ``firstfruits of repentance'' (Moroni 8:25; see \hyperref[FIRSTFRUITS]{Firstfruits}), and it is the way we witness to God that we are willing to repent, as stated in Alma 7:14--15:
				
					\begin{quote}
						Now I say unto you that ye must repent and be born again, for the Spirit saith: If ye are not born again, ye cannot inherit the kingdom of heaven. Therefore come and be baptized unto repentance, that ye may be washed from your sins, that ye may have faith on the Lamb of God, which taketh away the sins of the world, which is mighty to save and to cleanse from all unrighteousness. 15 Yea, I say unto you: Come and fear not, and lay aside every sin which easily doth beset you, which doth bind you down to destruction. Yea, come and go forth, and show unto your God that ye are willing to repent of your sins and enter into a covenant with him to keep his commandments, and witness it unto him this day by going into the waters of baptism.
					\end{quote}
					
				I think there is just something special about baptism. Clearly it's very symbolic---of Christ's death, the cleansing water, the waters of life, and other symbols I'm sure I could think of---but there also just needed to be a start to the path, and baptism happens to be it. 
				
				Although it's obviously important to keep all the commandments, or to do as much as possible to keep them, baptism is a very special one that, when fulfilled, brings a remission of sin.
				
				What about those of us who have already been baptized? Here I want to note a connection to Alma 32:16, which reads:
				
					\begin{quote}
						Therefore blessed are they who humbleth themselves without being compelled to be humble. Or rather, in other words, blessed is he that believeth in the word of God and is baptized without stubbornness of heart, yea, without being brought to know the word---or even compelled to know--- before they will believe. 
					\end{quote}
					
				After baptism, we are prompted by the Spirit or other people to make changes in our lives and enhance our level of righteousness. This will probably involve humbling ourselves in some way, and if we can do that, and be ``baptized without stubbornness of heart''---if we can fulfill whatever commandment we have been given, baptism representing here any possible instruction from the Lord---then we will be eligible for similar blessings, I think.
				
				The other point regarding those of us who have already been baptized regards the sacrament.
		
	% -----
	\subsubsection{The Lord's willingness to grant remission of sins}
	% -----
		\label{REMISSION-willingness}
		
		I don't have much to say about this yet, but in studying remission of sin I have been struck by how readily the Lord grants remissions (forgiveness) in the DC. It happens all the time to many people in many different circumstances.
		
		Some cases:
		
			\begin{compactitem}
				\item \hyperref[DC29-3]{DC 29:3}
				\item \hyperref[DC31-5]{DC 31:5}
				\item \hyperref[DC36-1]{DC 36:1}
				\item \hyperref[DC60-7]{DC 60:7}
				\item \hyperref[DC61-2]{DC 61:2}
				\item \hyperref[DC62-3]{DC 62:3}
				\item \hyperref[DC64-3]{DC 64:3}
				\item \hyperref[DC64-7]{DC 64:7}
				\item \hyperref[DC84-61]{DC 84:61}
				\item \hyperref[DC90-1]{DC 90:1}
				\item \hyperref[DC110-5]{DC 110:5}
				\item \hyperref[DC112-1]{DC 112:1--3}
				\item \hyperref[DC124-76]{DC 124:76, 78}
			\end{compactitem}
			
		Goes along with the Lord's willingness to heal in other places:
		
			\begin{compactitem}
				\item \hyperref[MATTHEW-8-2]{Matthew 8:2--3}
			\end{compactitem}
		
	% -----
	\subsubsection{Retaining a remission of sin}
	% -----
		\label{REMISSION-retaining}
		
		See especially \hyperref[MOSIAH-4-11]{Mosiah 4:11--12}, but there are other verses that talk about this idea.
		
		The key actions in Mosiah 4:11 are:
		
			\begin{compactitem}
				\item remember and always retain in remembrance the greatness of God and your own nothingness and his goodness and long-suffering towards you unworthy creatures
				\item humble yourselves even in the depths of humility
				\item calling on the name of the Lord daily
				\item standing steadfastly in the faith of that which is to come
			\end{compactitem}
		
	% -----
	\subsubsection{Remission and justification}
	% -----
		\label{REMISSION-justification}
		
		Remission and justification seem to have a lot in common---maybe it's valuable to equate them? Worth looking at in the future.